\chapter{}

Cheguei a Teresina com uma mala de roupas e três ou quatro livros.
Oficialmente, tinha sido chamada para fazer uma avaliação do andamento pedagógico do Colégio, por solicitação do novo reitor, Pe. Anchieta Lima Costa.
Essa era uma estratégia tipicamente jesuíta.
Eu devia circular pelos corredores, salas de aula e reuniões, observando, perguntando e anotando, sem dar a conhecer minha real condição de futura coordenadora-geral da escola, principalmente para as principais interessadas, as funcionárias que vinham comandando o trabalho pedagógico da escola.
Era uma situação que me punha muito desconfortável, até porque a farsa não enganou o pessoal mais antigo da casa que já me conhecia.
Alguns até se puseram a festejar abertamente o meu retorno.
Meu mal-estar só foi desaparecendo na medida em que a minha indignação crescia diante do que eu constatava.
Ao longo dos anos em que eu estivera ausente, aquela casa tinha sido invadida pelas práticas que por séculos vinham determinando o atraso e a pobreza daquela sociedade.
De agente de transformação, o Colégio tinha virado um santuário de preservação do que grassava de pior na cultura piauiense: empreguismo, clientelismo, submissão aos poderosos, descompromisso e irresponsabilidade.
Pedi um levantamento estatístico do aproveitamento dos alunos ao departamento de informática e levei-o ao Pe. Anchieta, juntamente com o resultado das minhas observações e anotações.
Era uma terça-feira, no final da tarde.
Pe. Anchieta viu, ouviu e determinou:

\textit{``-- Os professores têm semanalmente, nas quintas-feiras à tarde, uma reunião de duas horas para discussão e atualização do fazer pedagógico, comandada pelo Diretor.
Nessa próxima quinta-feira você assume a condução do encontro, já como a nova Coordenadora-Geral da escola.
Vou providenciar a imediata demissão das coordenadoras em exercício.''}

Eu estava hospedada na Casa dos Padres até que a imobiliária encontrasse um apartamento para me alugar.
Voltei para o meu quarto muito preocupada.
Tinha dois dias para me organizar.
Em primeiro lugar, assaltava-me a insegurança gerada pelo longo afastamento do ambiente escolar.
A rigor, eu estava melhor preparada para falar sobre peixes, panelas e vassouras do que sobre ensino e aprendizagem.
Em segundo, à exceção de uns poucos remanescentes da minha época, muitos daqueles professores faziam parte da ``República do Acará'', referência jocosa ao fato de residirem no mesmo bairro onde moravam as coordenadoras em vias de serem demitidas.
Num alegre cordão em que um vizinho indicava um conhecido, um parente puxava o outro, aquela gente toda tinha sido introduzida no Diocesano não por mérito ou capacidade, mas pela singela razão de que aquelas duas senhoras, muito relacionadas naquela região da cidade, achavam natural usar suas atribuições no colégio para ajeitar a situação de todos os parentes, amigos e conhecidos que estivessem necessitados de um emprego, quem sabe vitalício, assegurado pela complacência compulsória da Santa Madre Igreja.
Nesse clima de mútuos favores, é óbvio, não interessava a ninguém expor as falhas ou ineficiências desse ou daquele, para não perturbar a harmonia do ambiente.
Se os alunos não estavam aprendendo, o problema certamente era deles mesmos que não queriam nada com nada.
E, se aparecia algum pai descontente, acabava por se calar quando, ao final do ano, aquele filho condenado a uma reprovação certa, era miraculosamente guindado à série seguinte, a título de uma aposta cristã na sua capacidade de regeneração.

Essa maneira de apropriar-se em benefício próprio das instituições, é importante ressaltar, estava de tal forma enraizada nos costumes daquelas pessoas que elas não viam mal algum nisso, muito pelo contrário.
Dessa maneira, como já acontecera na minha experiência anterior no Colégio, eu já previa enfrentar muita indignação e hostilidade na medida em que começássemos a desfazer aquele tão cômodo, fraterno e deletério arranjo.

Na quinta-feira, munida dos levantamentos feitos pelos meninos do setor de informática, bastante estimulados pelo meu pedido de que os fizessem sob a forma de gráficos, os mais multicoloridos e impactantes que pudessem, apresentei-me diante do salão lotado por aquelas almas desassossegadas pelos boatos de demissões iminentes.
Eu também não estava lá muito firme sobre as pernas e, por isso mesmo, após um rápido ``boa tarde'', dei ordem ao auxiliar para projetar os gráficos sobre o telão e fiquei aguardando as reações, enquanto me preparava mentalmente para começar a espinhosa reunião.
Foi quando uma professora, considerada o destaque máximo do departamento de Português, segundo tinham me informado, apontou para o mais descendente dos desempenhos projetados, o de um primeiro colegial que ela lecionava, e exclamou, às gargalhadas:

\textit{``-- É isso aí.
E, no que depender de mim, vai cair muito mais no próximo bimestre.
Podem apostar!''}

Aquela intervenção foi bendita.
O sangue, fervendo, veio subindo até explodir no cérebro, como um ``estalo de Vieira'', um Pentecostes.
Não sei de onde veio, como veio, mas com voz pausada e controlada comecei a falar.
A sensação era de que eu estava tomada, porque as ideias me saíam claras, redondas, lógicas.
Fui me tornando mais e mais segura, mais incisiva e minhas convicções, para minha surpresa, foram aflorando formatadas, nítidas, a tal ponto que sentia estar sendo perfeitamente didática ao expô-las.
Falei, interpelei, desestabilizei, sacudi aquela plateia, mexi-lhes com as entranhas até fazê-los admitir que aquele fracasso, exposto em cores brilhantes na tela, \textbf{era sim a tradução irrefutável do trabalho profissional desenvolvido por eles com seus alunos}.
E que, para justificar aquele descalabro perante qualquer fórum minimamente sério, seria necessário bem mais que uma gargalhada irônica.
Em seguida, anunciei que, a partir daquela reunião, eu me dispunha a juntar-me a eles num trabalho árduo, sério e necessário de reabilitação pedagógica do Colégio Diocesano, contratada que eu fora pela nova diretoria para ocupar a Coordenação Geral da escola.

A tal professora teve um prolongado, e para mim bastante conveniente, ataque histérico.
Deixei-a gritar e lamentar-se até que a voz lhe faltou.
Acostumada a que os padres aflitos a acudissem com água de açúcar, pareceu chocada com minha impassibilidade e finalmente se calou quando se deu conta de que acabara de protagonizar um solitário espetáculo de absoluto descontrole, dando-me, ainda por cima, mais um excelente argumento para pedir seu desligamento da casa.

Refletindo depois sobre o que acontecera naquela tarde, concluí que o meu tempo ao longo dos últimos anos não fora afinal tão perdido como eu pensava.
Aprendera muito com a vida sobre as pessoas e, principalmente, ganhara algum terreno na caminhada pela superação dos meus receios, das minhas fragilidades, das minhas dúvidas sobre mim mesma.
Quantas vezes, ao tentar ler e me informar para não me sentir alijada do mundo, eu me entristecera ao descobrir que não retinha nada, fosse pelo cansaço, ou pela falta de uma aplicação imediata, porque esquecia completamente o que havia lido no dia anterior e, no dia seguinte, tinha que voltar ao ponto de partida.
Talvez, pensei, eu não tivesse guardado os ``o quês'' e os ``comos'', porém, acabava de descobrir, havia me aplicado ardentemente aos ``por quês'' e ``para quês'', decisão que me serviria muitíssimo daí por diante e cada vez mais, até os dias de hoje.

Poucos dias depois, a imobiliária conseguiu-me um pequeno apartamento, exatamente como eu queria.
Comprei de segunda mão uma cama, uma escrivaninha e algumas cadeiras.
Eu estava adorando aquela simplicidade monacal.
 Ela servia muito bem ao meu projeto de dedicar-me integralmente ao trabalho.
A última coisa que eu queria na vida, naquela altura, era me preocupar com casa, limpeza, cozinha.
Após incontáveis anos, eu tinha finalmente a oportunidade de ganhar meu próprio dinheiro, fazendo uma coisa importante e necessária e que eu começava a sentir que faria bem.
Não que o que eu fizera antes tivesse sido dispensável, sem importância ou menos gratificante.
Mas, se tudo caminha como o esperado, função de mãe é auto-extinguível, mesmo que o título seja vitalício.
E, com a partida dos meninos, tudo o que me restara era permanecer como uma espécie de curadora de um pequeno museu de lembranças familiares.
Nem eles, nem eu, merecíamos que eu tivesse esse futuro.
Quanto ao Paulo, minha suspeita era de que ele começara acreditar, de fato, que a família e eu, em particular, fôramos os responsáveis por ele não ter chegado até onde queria.
Uma temporada sozinho, eu esperava, haveria de ajudá-lo a colocar as coisas no enquadramento correto, como vinha acontecendo comigo.


Tenho que ser honesta e confessar que em minha primeira temporada à frente da Coordenação do Diocesano, eu me pautara pela preocupação de dar aos \textbf{meus filhos} a oportunidade de ter acesso a uma boa faculdade.
Apenas estendera esse objetivo a todos os outros alunos.
Houve até uma mãe, esposa de um deputado, que me parou um dia à porta para perguntar, bastante agastada:

``-- Será que você pensa que isto aqui é São Paulo? Meu filho não precisa se matar nem para entrar na faculdade e nem para conseguir um emprego em qualquer lugar desse Piauí.
Ele é filho de deputado, você esqueceu?'' 

Fui investigar o motivo da interpelação e descobri que o moleque estava em recuperação na maioria das disciplinas, abortando irremediavelmente as férias da família.

Nessa segunda passagem pela coordenação, porém, eu estava livre para me concentrar nas reais necessidades daquela gente.
Eu tentaria ``calçar os sapatos deles'', orientar-me pelas aspirações deles, porque do meu trabalho anterior pouco ou nada restara.
Ainda que aquela mãe estivesse errada ao negar ao filho o direito de impor-se ao mundo por seus próprios méritos, ela fora esperta em farejar qualquer coisa de alienígena nos rumos imprimidos à escola.
Eu aprendera uma importante lição: para ajudar alguém a ir a algum lugar, é preciso fundamentalmente que esse alguém \textbf{queira chegar lá}.
Do contrário, a empreitada não se sustenta.
É um fundamento inaciano: sem a adesão do \textbf{querer}, nenhuma ação se concretiza.

Nesse meu retorno ao Piauí, eu encontrara as coisas um pouco diferentes.
Um governador levado ao poder pelo Partido dos Trabalhadores e fartamente provido de verbas federais pelo governo Lula, realizava um intenso trabalho de propagandear o estado pelo restante do país como a última fronteira a ser explorada.
Diante disso, ficava a cada dia mais difícil para as velhas oligarquias dominantes mantê-lo na submissão quase feudal que eu conhecera da primeira vez.
Investidores de vários cantos do país e até multinacionais vinham se estabelecendo na região e, ainda que lentamente, mudando perspectivas, critérios e comportamentos.
Ao menos, ficava um pouco mais fácil falar em termos de educação para o século XXI.


\begin{figure}
\centering
\includegraphics[width=\linewidth]{41/colégio-modernizado-flacsi.jpg}
\caption*{O colégio modernizado. Fonte: Federação Latino-americana de Colégios Jesuítas (FLACSI)}
\end{figure}


No Diocesano, pelos próximos quatro anos, estaríamos sob a batuta hesitante de um reitor inteligente, culto, vaidoso e impregnado da proverbial preguiça baiana.
Um homem paradoxal, dividido entre a sensibilidade extremada que o fragilizava e um voluntarismo imperial que o levava a gestos de audácia temerária e a decisões autocráticas, capazes de desconcertar sobremaneira funcionários, professores e pais.
Sempre se soube que o sistema decisório na Companhia era imperial, de fato.
Mas, em Padre Anchieta, o que confundia e irritava seus interlocutores habituais era a instabilidade.
Ora ele parecia tímido e omisso, ora radical e autoritário.
Porém, a mim, deu-me todo o respaldo necessário para desenvolver meu trabalho e divertiu-se um bocado ajudando-me a fazer do Diocesano um colégio de ponta.
Era um orador competente e um estudioso devotado à Cristologia.
Apesar disso, e porque o entediavam as questões comezinhas da rotina da escola, costumava delegar a mim e por mais vezes do que seria conveniente, a tarefa de representar a ``voz do trono'' nas ocasiões em que se fazia apropriado abordar a espiritualidade inaciana.
A tal ponto que muita gente perguntava, sem levar em conta meu rosto maquilado e meus lábios sempre pintados em vermelho vivo, se eu era uma ``freira jesuíta``.
Ele só tomava o cuidado de ler com antecedência o que eu diria e, quando necessário, esclarecia um ou outro ponto impreciso ou obscuro.

\begin{figure}
\centering
\includegraphics[width=\linewidth]{41/+pe-anchieta-centenário-rescan.jpeg}
\caption*{Com Pe. Anchieta  no centenário do Colégio.}
\end{figure}

Não sei se me será reservado o privilégio de encontrar Ignácio de Loyola na eternidade.
Talvez seja melhor não.
Porque às vezes, no calor do entusiasmo, receio ter tomado liberdades indevidas com seu legado espiritual.

Nessa segunda passagem pelo Diocesano, falei por quantas juntas tinha.
Foi um trabalho exaustivo de convencimento dos professores, dos pais, dos alunos e até dos próprios jesuítas, às vezes.
Foram dezenas de palestras e apresentações nas quais eu me servi dos ensinamentos de Santo Ignácio para argumentar, defender, reivindicar, demonstrar a importância de preparar os alunos para o mundo informatizado e globalizado do Segundo Milênio.
Nossos jovens davam-se conta das perspectivas que vinham se abrindo para o Piauí, mas pareciam ter dificuldade em perceber de que maneira podiam e deviam inserir-se nelas.
E Santo Ignácio me foi muito útil.
Esse pequeno e tenaz guipuscoano do Renascimento sabia tudo sobre como preparar e fortalecer homens para enfrentar grandes mudanças

\begin{figure}
\centering
\includegraphics[width=\linewidth]{41/na-sala-rescan.jpg}
\caption*{Na minha sala de trabalho.}
\end{figure}

\begin{figure}
\centering
\includegraphics[width=\linewidth]{41/+alunos-rescan.jpg}
\caption*{Com os alunos, na Feira de Arte e Ciências.}
\end{figure}

\begin{figure}
\centering
\includegraphics[width=\linewidth]{41/equipe-rescan.jpg}
\caption*{A equipe de educadores do Colégio Diocesano em 2006.}
\end{figure}

Um dia, cerca de três anos depois do meu retorno ao Piauí, uma tia do Paulo, Cecília, irmã do meu sogro, telefonou-me de São Paulo.
Queria localizar Paulo com urgência.
Contou-me que Dr.
José Ignácio, prestes a receber uma importante herança que lhe coubera pela morte da irmã mais velha, Maria, caíra nas garras de um bando constituído por uma ex-namorada, agora octogenária, e seus dois suspeitíssimos filhos, um casal de vigaristas conhecidos na alta roda paulistana.
 Cecília acrescentou ainda que o irmão vinha demonstrando uma preocupante confusão mental e que se Paulo não voltasse imediatamente para assumir o comando da situação, ela receava pelo pior.
Localizei Paulo em Santarém e transmiti-lhe o recado.
Quase na sequência, ele embarcou para São Paulo e, lá chegando, constatou que a tia tinha toda razão de estar alarmada.
Voltou a Santarém apenas para encerrar seus negócios e, em seguida, a contragosto, obrigou-se a voltar às suas origens paulistanas.

Durante o tempo em que permaneci em Teresina, Paulo tinha ficado em Santarém e vinha visitar-me a cada três ou quatro semanas, para alívio dos Padres que não entendiam lá muito bem o jeito como funcionava aquele casamento.
Diga-se de passagem, não eram os únicos.
Nem os filhos e nem os parentes em geral entendiam também.
Mas eu estava achando bom e Paulo parecia apoiar meu entusiasmo pelo trabalho.
E ele, ao que parecia, gostava de Santarém e da vida que tinha lá.
Entretanto, a partir daquele telefonema, as coisas mudaram.
Ele ficou alguns meses em São Paulo, durante a fase aguda dos desvarios do pai.
Mas, como sempre odiou a cidade, assim que as coisas se acalmaram um pouco, decidiu se estabelecer em Cananéia.
Uma propriedade da família, na vizinha Ilha Comprida, que o pai deixara há muito de frequentar, estava sendo grilada e ele teve que iniciar um cansativo processo de expulsão judicial dos invasores.

Com Paulo de volta, nossos filhos começaram a me pressionar para ter a família de novo reunida.
Comecei a cogitar sobre aposentar-me.

Eu havia me comprometido a permanecer em Teresina enquanto durasse a gestão de Padre Anchieta.
E ele chegaria ao término do mandato em poucos meses.
Eu já sabia que havia interesse da Companhia em levá-lo dali, embora ele lutasse por permanecer.
Mas, tudo se encaminhou para a sua partida e logo já se sabia até quem viria para o seu lugar.
Era Padre Darly, meu antigo chefe, com quem eu me dava bem, mas que nem de longe tinha os horizontes largos de Padre Anchieta.
De mais a mais, o que eu me propusera a fazer estava feito.
Era mais do que hora de entregar a escola a uma liderança local.

É minha convicção de que há um momento, para quem atua em Educação, em que é preciso parar.
É quando você ainda é capaz de entender o que se passa, mas já não é capaz de aceitar.
Mesmo em Teresina, uma sociedade conservadora, já eram sensíveis os efeitos das mudanças que vinham afetando a estrutura das famílias.
Metade do tempo dedicado a palestras naqueles quatro anos, fora gasto na tentativa de ajudar filhos e maridos a se adaptarem à ausência das mães, crescentemente empenhadas na batalha diária por suas carreiras profissionais.
Ou tentando corrigir, nas crianças, os efeitos de uma convivência doméstica orientada por valores esquálidos, consumismo desvairado e intolerância coletiva às frustrações.
Ou ainda, convencendo pais imaturos a exercerem sua autoridade, pelo bem da segurança emocional dos filhos.
A questão era que ou a escola assumia esse papel, ou as consequências de todos esses desajustes viriam desaguar no baixo desempenho dos alunos frente aos desafios do aprendizado.
Uma tendência parecia se configurar, e isso era confirmado em encontros com profissionais de outros colégios e de outros lugares, de as escolas irem assumindo progressivamente responsabilidades antes exclusivas das famílias, até porque, em muitos casos, constituíam o único ambiente organizado e estável na vida daqueles jovens.
Era uma perspectiva que não me motivava, basicamente porque eu acredito que família e escola podem e devem ser complementares na sua atuação, mas não vejo como uma possa suprir o papel da outra.
As reformas que implantáramos no Diocesano habilitaram-no a cumprir com excelência o papel de preparar aquela juventude para o exercício de uma cidadania produtiva, ética e participativa.
Quanto ao que a sociedade parece esperar para além disso, pareceu-me melhor entregar a questão a cabeças mais afinadas com os novos tempos.
Após acertar-me com Padre Anchieta, pedi minha aposentadoria, aos trinta e quatro anos de exercício profissional.

Na nossa despedida, em meio ao carinho e as homenagens que ele e eu recebemos, tive a enorme surpresa de deparar-me com Paula e Marcelo que, a convite dos padres e sem que eu sequer suspeitasse, chegaram a Teresina na companhia do pai para participar da festa.

\begin{figure}
\centering
\includegraphics[width=0.9\linewidth]{41/despedida-1-rescan-edit.jpg}
\end{figure}

\begin{figure}
\centering
\includegraphics[width=0.9\linewidth]{41/despedida-2-rescan-edit.jpg}
\caption*{Paula, Marcelo e Paulo chegando de surpresa para a minha despedida do Diocesano.}
\end{figure}

\begin{figure}
\centering
\includegraphics[width=\linewidth]{41/festa-despedida-rescan-edit.jpg}
\caption*{A festa da despedida, ao lado de Pe. Darly, Pe. Anchieta, Áurea e Albino.}
\end{figure}